%%%%%%%%%%%%%%%%%%%%%%%%%%%%%%%%%%%%%%%%%%%%%%%%%%%%%%%%%%%%%%%%%%%%%%
% Lớp 11 | Hàm số lượng giác và phương trình lượng giác | [Giảm] Phương trình lượng giác thường gặp
% Dạng 1D1?6-7: Phương trình lượng giác có chứa tham số
%%%%%%%%%%%%%%%%%%%%%%%%%%%%%%%%%%%%%%%%%%%%%%%%%%%%%%%%%%%%%%%%%%%%%%


% nguồn: 2026-11-15_11_1D1_TL2_Luong_giac (nhập 15/11/2026)
\begin{bt}%[1D1V6-7]%[ma:92108107]
% nguon: Archive/SP_KhoLatex/ID6_old/1K1-4-9.tex #2986
[VDT]%[DCHT Toán 11 - KNTT -Tên Nguyễn Thanh Sang] 
Tìm tất cả các giá trị nguyên của $m$ để phương trình $(m-1)\sin 2x-4\cos 2x=2-2m$ có nghiệm.
	\dapso{$-1, 0, 1, 2, 3.$}
\loigiai{
		Để phương trình có nghiệm thì \begin{eqnarray*}
		& &(m-1)^2+(-4)^2\ge (2-2m)^2\\
		&\Leftrightarrow& m^2-2m+1+16\ge 4-8m+4m^2\\
		&\Leftrightarrow& 3m^2-6m-13\le 0\\
		&\Leftrightarrow& \dfrac{3-4\sqrt{3}}{3}\le m\le \dfrac{3+4\sqrt{3}}{3}.
	\end{eqnarray*}
	Khi đó các giá trị nguyên của tham số $m$ là $-1$, $0$, $1$, $2$, $3$.
	}
\end{bt}

% nguồn: 2026-11-15_11_1D1_TL2_Luong_giac (nhập 15/11/2026)
\begin{bt}%[1D1V6-7]%[ma:31d7c2b5]
% nguon: Archive/SP_KhoLatex/ID6_old/1K1-4-9.tex #2987
[VDT]%[DCHT Toán 11 - KNTT -Tên Nguyễn Thanh Sang] 
	Tính tổng tất cả các giá trị nguyên của tham số $m$ để phương trình $(2m+1)\sin x-(m+2)\cos x=2m+3$ vô nghiệm.
	\dapso{$10.$}
	\loigiai{
	Phương trình vô nghiệm khi
	\begin{eqnarray*}
		& & (2m+1)^2 + (m+2)^2 < (2m+3)^2\\
		&\Leftrightarrow & 4m^2+4m+1 +m^2+4m+4<4m^2+12m+9\\
		&\Leftrightarrow & m^2-4m-4<0\\
		&\Leftrightarrow & 2-2\sqrt{2} < m < 2+2\sqrt{2}.
	\end{eqnarray*}
	Do $m$ nguyên nên ta nhận $m \in \{0;1;2;3;4\}$.\\
	Vậy tổng tất cả các giá trị nguyên của tham số $m$ bằng: $0+1+2+3+4=10$.
}
\end{bt}

% nguồn: 2026-11-15_11_1D1_TL2_Luong_giac (nhập 15/11/2026)
\begin{bt}%[1D1V6-7]%[ma:bd03c565]
% nguon: Archive/SP_KhoLatex/ID6_old/1K1-4-9.tex #2993
[VDT]%[DCHT Toán 11 - KNTT -Tên Nguyễn Thanh Sang]  
Tìm tất cả các giá trị của tham số m để phương trình $\sin x+\left(m-1\right)\cos x=2m-1$ có nghiệm.
	\dapso{$-\dfrac{1}{3}\le m\le 1.$}
	\loigiai{
	Phương trình có nghiệm khi: $1^2+(m-1)^2\ge (2m-1)^2 \Leftrightarrow 3m^2-2m-1\le 0 \Leftrightarrow -\dfrac{1}{3}\le m\le 1$.
		}
\end{bt}

% nguồn: 2026-11-15_11_1D1_TL2_Luong_giac (nhập 15/11/2026)
\begin{bt}%[1D1V6-7]%[ma:360c069c]
% nguon: Archive/SP_KhoLatex/ID6_old/1K1-4-9.tex #2994
[VDT]%[DCHT Toán 11 - KNTT -Tên Nguyễn Thanh Sang] 
	Tìm tất cả các giá trị thực của tham số $m$ để phương trình $m\sin 2x+(m-1)\cos 2x=\sqrt{5}m$ vô nghiệm.
	\dapso{$m>\dfrac{1}{3}$ hoặc $m<-1$.}
	\loigiai{
	Phương trình đã cho vô nghiệm khi và chỉ khi:
	\begin{align*}
		&m^2+(m-1)^2<\left(\sqrt{5}m\right)^2 \Leftrightarrow 2m^2-2m+1< 5m^2\\
		\Leftrightarrow\ &3m^2+2m-1>0 \Leftrightarrow \hoac{&m>\dfrac{1}{3}\\ &m<-1.}
	\end{align*}
	}
\end{bt}
